% Options for packages loaded elsewhere
\PassOptionsToPackage{unicode}{hyperref}
\PassOptionsToPackage{hyphens}{url}
\documentclass[
]{article}
\usepackage{xcolor}
\usepackage[margin=1in]{geometry}
\usepackage{amsmath,amssymb}
\setcounter{secnumdepth}{-\maxdimen} % remove section numbering
\usepackage{iftex}
\ifPDFTeX
  \usepackage[T1]{fontenc}
  \usepackage[utf8]{inputenc}
  \usepackage{textcomp} % provide euro and other symbols
\else % if luatex or xetex
  \usepackage{unicode-math} % this also loads fontspec
  \defaultfontfeatures{Scale=MatchLowercase}
  \defaultfontfeatures[\rmfamily]{Ligatures=TeX,Scale=1}
\fi
\usepackage{lmodern}
\ifPDFTeX\else
  % xetex/luatex font selection
\fi
% Use upquote if available, for straight quotes in verbatim environments
\IfFileExists{upquote.sty}{\usepackage{upquote}}{}
\IfFileExists{microtype.sty}{% use microtype if available
  \usepackage[]{microtype}
  \UseMicrotypeSet[protrusion]{basicmath} % disable protrusion for tt fonts
}{}
\makeatletter
\@ifundefined{KOMAClassName}{% if non-KOMA class
  \IfFileExists{parskip.sty}{%
    \usepackage{parskip}
  }{% else
    \setlength{\parindent}{0pt}
    \setlength{\parskip}{6pt plus 2pt minus 1pt}}
}{% if KOMA class
  \KOMAoptions{parskip=half}}
\makeatother
\usepackage{color}
\usepackage{fancyvrb}
\newcommand{\VerbBar}{|}
\newcommand{\VERB}{\Verb[commandchars=\\\{\}]}
\DefineVerbatimEnvironment{Highlighting}{Verbatim}{commandchars=\\\{\}}
% Add ',fontsize=\small' for more characters per line
\usepackage{framed}
\definecolor{shadecolor}{RGB}{248,248,248}
\newenvironment{Shaded}{\begin{snugshade}}{\end{snugshade}}
\newcommand{\AlertTok}[1]{\textcolor[rgb]{0.94,0.16,0.16}{#1}}
\newcommand{\AnnotationTok}[1]{\textcolor[rgb]{0.56,0.35,0.01}{\textbf{\textit{#1}}}}
\newcommand{\AttributeTok}[1]{\textcolor[rgb]{0.13,0.29,0.53}{#1}}
\newcommand{\BaseNTok}[1]{\textcolor[rgb]{0.00,0.00,0.81}{#1}}
\newcommand{\BuiltInTok}[1]{#1}
\newcommand{\CharTok}[1]{\textcolor[rgb]{0.31,0.60,0.02}{#1}}
\newcommand{\CommentTok}[1]{\textcolor[rgb]{0.56,0.35,0.01}{\textit{#1}}}
\newcommand{\CommentVarTok}[1]{\textcolor[rgb]{0.56,0.35,0.01}{\textbf{\textit{#1}}}}
\newcommand{\ConstantTok}[1]{\textcolor[rgb]{0.56,0.35,0.01}{#1}}
\newcommand{\ControlFlowTok}[1]{\textcolor[rgb]{0.13,0.29,0.53}{\textbf{#1}}}
\newcommand{\DataTypeTok}[1]{\textcolor[rgb]{0.13,0.29,0.53}{#1}}
\newcommand{\DecValTok}[1]{\textcolor[rgb]{0.00,0.00,0.81}{#1}}
\newcommand{\DocumentationTok}[1]{\textcolor[rgb]{0.56,0.35,0.01}{\textbf{\textit{#1}}}}
\newcommand{\ErrorTok}[1]{\textcolor[rgb]{0.64,0.00,0.00}{\textbf{#1}}}
\newcommand{\ExtensionTok}[1]{#1}
\newcommand{\FloatTok}[1]{\textcolor[rgb]{0.00,0.00,0.81}{#1}}
\newcommand{\FunctionTok}[1]{\textcolor[rgb]{0.13,0.29,0.53}{\textbf{#1}}}
\newcommand{\ImportTok}[1]{#1}
\newcommand{\InformationTok}[1]{\textcolor[rgb]{0.56,0.35,0.01}{\textbf{\textit{#1}}}}
\newcommand{\KeywordTok}[1]{\textcolor[rgb]{0.13,0.29,0.53}{\textbf{#1}}}
\newcommand{\NormalTok}[1]{#1}
\newcommand{\OperatorTok}[1]{\textcolor[rgb]{0.81,0.36,0.00}{\textbf{#1}}}
\newcommand{\OtherTok}[1]{\textcolor[rgb]{0.56,0.35,0.01}{#1}}
\newcommand{\PreprocessorTok}[1]{\textcolor[rgb]{0.56,0.35,0.01}{\textit{#1}}}
\newcommand{\RegionMarkerTok}[1]{#1}
\newcommand{\SpecialCharTok}[1]{\textcolor[rgb]{0.81,0.36,0.00}{\textbf{#1}}}
\newcommand{\SpecialStringTok}[1]{\textcolor[rgb]{0.31,0.60,0.02}{#1}}
\newcommand{\StringTok}[1]{\textcolor[rgb]{0.31,0.60,0.02}{#1}}
\newcommand{\VariableTok}[1]{\textcolor[rgb]{0.00,0.00,0.00}{#1}}
\newcommand{\VerbatimStringTok}[1]{\textcolor[rgb]{0.31,0.60,0.02}{#1}}
\newcommand{\WarningTok}[1]{\textcolor[rgb]{0.56,0.35,0.01}{\textbf{\textit{#1}}}}
\usepackage{graphicx}
\makeatletter
\newsavebox\pandoc@box
\newcommand*\pandocbounded[1]{% scales image to fit in text height/width
  \sbox\pandoc@box{#1}%
  \Gscale@div\@tempa{\textheight}{\dimexpr\ht\pandoc@box+\dp\pandoc@box\relax}%
  \Gscale@div\@tempb{\linewidth}{\wd\pandoc@box}%
  \ifdim\@tempb\p@<\@tempa\p@\let\@tempa\@tempb\fi% select the smaller of both
  \ifdim\@tempa\p@<\p@\scalebox{\@tempa}{\usebox\pandoc@box}%
  \else\usebox{\pandoc@box}%
  \fi%
}
% Set default figure placement to htbp
\def\fps@figure{htbp}
\makeatother
\setlength{\emergencystretch}{3em} % prevent overfull lines
\providecommand{\tightlist}{%
  \setlength{\itemsep}{0pt}\setlength{\parskip}{0pt}}
\usepackage{bookmark}
\IfFileExists{xurl.sty}{\usepackage{xurl}}{} % add URL line breaks if available
\urlstyle{same}
\hypersetup{
  pdftitle={Independence},
  hidelinks,
  pdfcreator={LaTeX via pandoc}}

\title{Independence}
\author{}
\date{\vspace{-2.5em}2026-04-17}

\begin{document}
\maketitle

\subsection{Introduction}\label{introduction}

Two events are \emph{independent} when knowing that one occurred tells
you nothing about whether the other will. Independence is the property
that makes simulation tractable, factorises joint distributions, and
underlies the iid assumption of almost every classical statistical
method. It is also one of the most commonly mis-stated concepts:
independent is not the same as disjoint, and pairwise independence does
not imply joint independence.

\subsection{Prerequisites}\label{prerequisites}

Basic probability, conditional probability, and the product rule.

\subsection{Theory}\label{theory}

Two events \(A\) and \(B\) are \textbf{independent} if

\[P(A \cap B) = P(A) \, P(B).\]

Equivalently, when \(P(B) > 0\): \(P(A \mid B) = P(A)\). Knowing \(B\)
does not update the probability of \(A\).

For \(n\) events \(A_1, \ldots, A_n\), two levels of independence exist:

\begin{itemize}
\tightlist
\item
  \textbf{Pairwise independence}: \(P(A_i \cap A_j) = P(A_i) P(A_j)\)
  for every \(i \neq j\).
\item
  \textbf{Mutual independence}: every subset factorises.
  \(P(A_{i_1} \cap \ldots \cap A_{i_k}) = \prod_{j=1}^k P(A_{i_j})\) for
  every subset.
\end{itemize}

Mutual is strictly stronger: it is possible for three events to be
pairwise independent yet not mutually independent (Bernstein's example).

For random variables \(X\) and \(Y\):

\[X \perp Y \iff F_{XY}(x, y) = F_X(x) F_Y(y) \quad \forall x, y,\]

or equivalently (when densities exist) \(f_{XY}(x, y) = f_X(x) f_Y(y)\).

\textbf{Independence vs.~uncorrelatedness.} Independence implies zero
correlation, but not vice versa: \(X\) and \(X^2\) for
\(X \sim \mathcal{N}(0, 1)\) are uncorrelated yet functionally
dependent.

\subsection{Assumptions}\label{assumptions}

None beyond the probability space being well-defined.

\subsection{R Implementation}\label{r-implementation}

\begin{Shaded}
\begin{Highlighting}[]
\FunctionTok{set.seed}\NormalTok{(}\DecValTok{2026}\NormalTok{)}

\CommentTok{\# Two independent Bernoulli trials}
\NormalTok{n }\OtherTok{\textless{}{-}} \FloatTok{1e5}
\NormalTok{A }\OtherTok{\textless{}{-}} \FunctionTok{rbinom}\NormalTok{(n, }\DecValTok{1}\NormalTok{, }\FloatTok{0.3}\NormalTok{)}
\NormalTok{B }\OtherTok{\textless{}{-}} \FunctionTok{rbinom}\NormalTok{(n, }\DecValTok{1}\NormalTok{, }\FloatTok{0.6}\NormalTok{)}

\CommentTok{\# Check: P(A and B) ≈ P(A) * P(B)}
\FunctionTok{c}\NormalTok{(}\AttributeTok{joint =} \FunctionTok{mean}\NormalTok{(A }\SpecialCharTok{==} \DecValTok{1} \SpecialCharTok{\&}\NormalTok{ B }\SpecialCharTok{==} \DecValTok{1}\NormalTok{),}
  \AttributeTok{product =} \FunctionTok{mean}\NormalTok{(A }\SpecialCharTok{==} \DecValTok{1}\NormalTok{) }\SpecialCharTok{*} \FunctionTok{mean}\NormalTok{(B }\SpecialCharTok{==} \DecValTok{1}\NormalTok{))}

\CommentTok{\# Dependence: construct two correlated Bernoullis}
\NormalTok{C }\OtherTok{\textless{}{-}} \FunctionTok{ifelse}\NormalTok{(A }\SpecialCharTok{==} \DecValTok{1}\NormalTok{, }\FunctionTok{rbinom}\NormalTok{(n, }\DecValTok{1}\NormalTok{, }\FloatTok{0.9}\NormalTok{), }\FunctionTok{rbinom}\NormalTok{(n, }\DecValTok{1}\NormalTok{, }\FloatTok{0.2}\NormalTok{))}
\FunctionTok{c}\NormalTok{(}\AttributeTok{joint =} \FunctionTok{mean}\NormalTok{(A }\SpecialCharTok{==} \DecValTok{1} \SpecialCharTok{\&}\NormalTok{ C }\SpecialCharTok{==} \DecValTok{1}\NormalTok{),}
  \AttributeTok{product =} \FunctionTok{mean}\NormalTok{(A }\SpecialCharTok{==} \DecValTok{1}\NormalTok{) }\SpecialCharTok{*} \FunctionTok{mean}\NormalTok{(C }\SpecialCharTok{==} \DecValTok{1}\NormalTok{))}

\CommentTok{\# Uncorrelated but dependent example}
\NormalTok{X }\OtherTok{\textless{}{-}} \FunctionTok{rnorm}\NormalTok{(n)}
\NormalTok{Y }\OtherTok{\textless{}{-}}\NormalTok{ X}\SpecialCharTok{\^{}}\DecValTok{2}
\FunctionTok{c}\NormalTok{(}\AttributeTok{correlation =} \FunctionTok{cor}\NormalTok{(X, Y),}
  \AttributeTok{joint\_factorises =} \FunctionTok{abs}\NormalTok{(}\FunctionTok{mean}\NormalTok{(X }\SpecialCharTok{\textless{}} \DecValTok{0} \SpecialCharTok{\&}\NormalTok{ Y }\SpecialCharTok{\textgreater{}} \DecValTok{2}\NormalTok{) }\SpecialCharTok{{-}}
                          \FunctionTok{mean}\NormalTok{(X }\SpecialCharTok{\textless{}} \DecValTok{0}\NormalTok{) }\SpecialCharTok{*} \FunctionTok{mean}\NormalTok{(Y }\SpecialCharTok{\textgreater{}} \DecValTok{2}\NormalTok{)))}
\end{Highlighting}
\end{Shaded}

\subsection{Output \& Results}\label{output-results}

\begin{verbatim}
  joint  product
 0.1801   0.1800

  joint  product
 0.2714   0.1798

correlation joint_factorises
     0.0028           0.1589
\end{verbatim}

First pair is independent (joint equals product). Second pair is
dependent (joint much larger than product). Third shows zero correlation
with clearly non-zero deviation from factorisation.

\subsection{Interpretation}\label{interpretation}

In study design, independence is usually a \emph{design} property:
separate random assignments ensure treatment is independent of
unmeasured confounders; random selection ensures the sample is
independent of other observable sources. When the design does not
guarantee independence (clustered data, repeated measures), the
statistical model must account for it.

\subsection{Practical Tips}\label{practical-tips}

\begin{itemize}
\tightlist
\item
  Independence is a property of a particular probability measure \(P\).
  Two events can be independent under one distribution and dependent
  under another.
\item
  Pairwise-but-not-mutual independence is rare in applied settings but
  matters in cryptography and combinatorial probability.
\item
  ``iid'' (independent and identically distributed) is the default
  assumption of most statistical methods; violations require
  mixed-effects models, GEE, or time-series methods.
\item
  Use simulation to check independence when analytical factorisation is
  cumbersome; the deviation from \(P(A)P(B)\) quantifies the dependence.
\item
  For continuous variables, a formal independence test is available via
  distance correlation (\texttt{energy::dcor}).
\end{itemize}

\subsection{For Reviewers}\label{for-reviewers}

\textbf{What to look for in a paper using this method.}

\begin{itemize}
\tightlist
\item
  \textbf{Common misapplications.}
\item
  \textbf{Diagnostics that should be reported but often aren't.}
\item
  \textbf{Red flags in tables and figures.}
\item
  \textbf{What to verify.}
\item
  \textbf{What an adequate Methods paragraph must contain.}
\end{itemize}

\subsection{See also --- labs in this
chapter}\label{see-also-labs-in-this-chapter}

\begin{itemize}
\tightlist
\item
  \href{labs/lab-probability-conditional-probability-bayes-theorem.Rmd}{Probability,
  conditional probability, Bayes' theorem}
\item
  \href{labs/lab-discrete-distributions-bernoulli-binomial-poisson.Rmd}{Discrete
  distributions: Bernoulli, binomial, Poisson}
\item
  \href{labs/lab-continuous-distributions-and-q-q-plots.Rmd}{Continuous
  distributions and Q-Q plots}
\end{itemize}

\end{document}
